\subsection{Pure links and simultaneous pruning}
\label{fa:links}
We prove \Cref{r:forward-marker} by successive exclusions, always stopping if the required point
stabilizers have already appeared.  The present subsection records the transfer and selection facts that
allow every exclusion to be made simultaneously below each bounded path prefix.  All labels refer to the
original marker set of \Cref{r:pure-markers}; ordinary labels are those represented by products at the
original simple sites.

\begin{lemma}[Exclusion of bounded types]\label{fa:bounded-types}
In the product configuration of \Cref{r:forward-marker}, the ambient simple group has unbounded order and
unbounded classical rank or alternating degree as the required marker reservoir grows.  More precisely,
any prescribed bounds on order and on Lie rank can be excluded by choosing the reservoir sufficiently
large.  This assertion also holds in every exact pure-marker filter obtained by the changes of simple
sites below.
\end{lemma}

\begin{proof}

Bounded order is immediately impossible by the length of an axis chain. In bounded Lie rank use the
argument of \Cref{w:lp,w:height-rigidity,w:bounded-exclusion}. For a prescribed sufficiently large constant \(k\) test \(M_U\),
\(|U|=k\). There must be some for which all ordinary labels above it admit the sandwich on intersection
with \(S\): otherwise pick many disjoint such sets and take proper bounded-equation stabilizer witnesses
as there. Each has omitted (nonabsorbed) set nonempty inside its tested set, hence intersections strictly
descend (also ordinary products, strict on intersection with \(S\)). The linear spaces of equations
cannot give such long descent. Now residual lengths under \(M_U\) are arbitrarily large and the same
height and order comparison for the sandwich maps bounds the number of ordinary labels above it by
\(C2^{h}\), \(h=O(k\log\log(k+3))\). But all partitions of \(U\) occur via intersections of cuts \(C_D\)
(individual symmetric group on every part), giving at least \(\exp((1/2+o(1))k\log k)\) for example by
pairings. This is impossible. In particular we can require arbitrarily high thresholds on the remaining
types here or at any use with a large enough marker set.

\end{proof}

\begin{convention}[Relative labels at an ancestor]\label{fa:relative-labels}
At a path prefix \(A\), suppose that we have the prepared-path data of
\Cref{w:prepared-path}: an actual simple subgroup
\[
 B_A\trianglelefteq K_{M_A},\qquad B_A\le S,\qquad HB_A=K_{M_A},
\]
whose centralizer in \(K_{M_A}\) is soluble and contained in \(H\).  Initially \(A=\varnothing\).  We
work in its almost simple quotient, and write \(T'=T\setminus A\) for the relative marker set.  Thus the
relative top is \(M_A\), even when a displayed relative label uses only \(T'\).
A proper full stabilizer containing a tested axis with a bounded number of further omissions has a product label
\begin{equation}\label{fa:witness-labels}
 r_D=(T'\setminus D,V_D),\qquad
 g_D=(\varnothing,V_D)=r_D\wedge X_D,
 \qquad \varnothing\ne D\subseteq T',\quad V_D\le S_D.
\end{equation}
Here \(D\) is a subset of those further omissions, and \(V_D\) fixes every position outside \(D\) in the
expression for \(g_D\).  A normal subgroup is \emph{retained} at this link if it is contained in
\(K_{g_D}\); otherwise it is \emph{missing}.
\end{convention}

We will use exact pure-marker links before passing to the ordinary site
interval.  These are different intervals: the ordinary one corresponds
exactly to the coordinate-product labels in the full formal link, and
need not contain every formal label.  Normal subgroups can nevertheless
be tested in the full link by transporting them across the original
simple sites.  We first prove this normal-subgroup transfer, then use
it to find independently supported groups inside intersections.

\begin{lemma}[Normal transfer on pure links]\label{fa:pure-links}
For the labels in \Cref{fa:witness-labels}, the full formal interval \([g_D,r_D]\) is \(F_{T'\setminus
D}\).  If this residual set is sufficiently large, its ordinary group interval \([K_{g_D},K_{r_D}]\) has
the normal property of \Cref{r:marker-normal}, including its nonabelian monolith modulo the base core.
For pairwise disjoint supports \(D_i\), put
\[
 J=\bigwedge_i r_{D_i},\qquad g_* =\bigvee_i g_{D_i},\qquad D_* =\bigcup_i D_i.
\]
The same conclusion holds for \([K_{g_*},K_J]\), and for distinct \(i,j\) one has \(g_{D_i}\le r_{D_j}\)
and \(r_{D_i}\vee r_{D_j}=M_A\).
\end{lemma}

\begin{proof}

The full formal link \([g_D,r_D]\) is \(F_{T'\setminus D}\):
every intermediate increase keeps \(V_D\) independent.  For disjoint
supports the label calculation gives \(g_{D_i}\le r_{D_j}\) and
\(r_{D_i}\vee r_{D_j}=M_A\).  The link \([g_*,J]\) is likewise
pure on \(T'\setminus D_*\), with independent side group
\(\prod_i V_{D_i}\).  We prove the normal assertion for either
pair of endpoints, denoted by \(g,r\).

Work before any ancestor-centralizer quotient, and let \(S\) denote
the original simple site.  Put
\[
 L=K_g,\qquad R=K_r,\qquad Q=L\cap S,\qquad Z=R\cap S.
\]
Let \(P_g,P_r\) be the corresponding products at the original sites,
with fixed-site projections \(Q,Z\).  They have these projections
because \(\Lambda\cap S\) lies in the original base.  With
\(\Delta\) replaced by its finite acting image if necessary, the
full formal link has the exact ordinary realization
\[
 [P_g\rtimes\Delta,\ P_r\rtimes\Delta].
\]
Indeed its intermediates all contain \(\Delta\), and hence correspond
to precisely the \(\Delta\)-invariant subgroups between the two
products.  The normal test applies to this full interval.

First let \(N\trianglelefteq R\) lie in \(Z\).  Since
\(\Lambda\le L\le R\), its transport to the other sites is
well-defined.  The resulting product is normal in
\(P_r\rtimes\Delta\).  By \Cref{r:marker-normal}, either
\(N\le Q\) or \(QN=Z\); in the second case \(LN=R\).
Thus every missing normal inside the original site supplements the
ordinary lower endpoint.

Now let \(N\trianglelefteq R\) satisfy \(N\cap Z\le L\).
Then \([N,Z]\le Q\), so \(N\) normalizes the intersection with
\(S\) of each singleton-omission axis coatom of the link.  These
coatoms are products, are strictly intermediate, and have meet \(L\).
If \(N\) escaped one of them, adjoining it to that coatom would
generate \(R\); the coatom's inner intersection would therefore be
normal in \(R\).  It strictly contains \(Q\), but its product
with \(L\) is only the proper coatom, contradicting the preceding
paragraph.  Hence \(N\le L\).  Together the two paragraphs prove
the normal endpoint assertion for every normal subgroup of \(R\).

Finally put \(C=\core_R(L)\).  Two missing normals commuting modulo
\(C\), or one missing normal abelian modulo \(C\), would transport
their intersections with \(Z\) to missing normal products in the
exact semidirect interval.  Their commutators lie in the transport of
\(C\cap Z\), a normal product contained in its base.  This contradicts
the unique nonabelian monolith supplied there by \Cref{r:marker-normal}.
The ordinary interval therefore has property~\((*)\).  Quotienting
an ancestor centralizer preserves it, since that kernel already lies
in \(L\).

\end{proof}

The lower endpoint for retention is thus \(K_{g_D}\) on the link
indexed by \(D\), and \(K_{g_*}\) on a simultaneous intersection link.
We keep these endpoints when applying the subsequent normal-product
arguments.

\begin{lemma}[The disjoint-support join restriction]\label{fa:disjoint-join}
For disjoint nonempty \(D_1,D_2\) and \(J=r_{D_1}\wedge r_{D_2}\), every formal label \(u\) satisfies
\begin{equation}\label{fa:join-restriction}
 J\le u<r_{D_1}\quad\Longrightarrow\quad u\vee r_{D_2}<M_A.
\end{equation}
\end{lemma}

\begin{proof}

Indeed, all positions outside \(D_1\cup D_2\) are already absorbed at \(u\). The part of \(D_2\) still
not absorbed there stays separate also from \(D_1\) since \(u\le r_{D_1}\), and is preserved by
\(V_{D_2}\) since \(J\le u\); if this part were empty one would already have \(r_{D_1}\).

\end{proof}

\begin{lemma}[Independent support in intersections]\label{fa:intersection-support}
In the setting of \Cref{fa:pure-links}, suppose that \(Y\trianglelefteq K_{r_D}\), with \(Y\le B_A\), and
that a retained normal subgroup \(Z\le Y\cap K_{g_D}\) makes \(Y/Z\) a missing power of nonabelian simple
groups.  Assume that \(K_{g_D}\) acts transitively on its simple sites with a bounded number of orbits on
ordered pairs.  At a sufficiently long residual link, singleton axis intersections with this power are
products and their simple sites have arbitrarily high type.  Moreover, in \(E=K_J\), the product of the
individual simple-site intersections of \((E\cap Y)/Z\) is \(E\)-normal and missing relative to
\(K_{g_*}\).  It has nonsoluble factors and its image contains the nonabelian monolith modulo the base
core of \(E\).  If \(J<r_D\), these singly supported factors are proper.
\end{lemma}

\begin{proof}

There are two sets of coordinates in this argument: the original \(b\)
simple sites, and the simple factors of \(Y/Z\) inside each of them.
We first obtain an exact interval at the new sites, then show that its
axis groups contain independent supports.

Lift \(Y,Z\) inside the actual \(B_A\le S\), before any ancestor
quotient, and transport them across the original sites.  Since
\(K_{g_D}Y=K_{r_D}\), we also have
\[
 (K_{g_D}\cap S)Y=K_{r_D}\cap S.
\]
Thus the transported product of copies of \(Y\) supplements the full
lower endpoint, while the transported product of copies of \(Z\)
is retained.  Every intermediate subgroup is its lower endpoint
times its intersection with that normal product.  After factoring
the retained product, \([g_D,r_D]\) is therefore an exact
invariant-subgroup interval in \((Y/Z)^b\).  The operators are the
original action together with conjugation by the label at \(g_D\);
they act transitively on the new simple sites.

Consider a singleton axis \(g_D\vee X_{\{i\}}\), with \(i\) outside
the used supports.  Its intersection with the new power remains a
product across the original \(b\) sites.  It has nontrivial projections
at the new sites, by strict increase from the lower endpoint and
transitivity.  These projections cannot be full.  Otherwise every
subgroup on a long axis chain above it would be subdirect.  Strict
axis increases give strict intersection increases by supplementation;
inside an original site these would require strict refinements of
\(K_{g_D}\)-invariant tie partitions.  The bounded number of orbits
on ordered pairs bounds the length of such a chain, as in
\Cref{c:orbit-bound}, whereas the residual marker set is arbitrarily
large.

The filter above this singleton axis is again an exact pure-marker
interval.  Its bottom has nontrivial proper projections to the new
sites, so \Cref{r:marker-mixed} makes it a coordinate product there.
\Cref{r:marker-site} then makes all subsequent axes products at those
sites as well.  Applying \Cref{fa:bounded-types} to this exact filter
forces its new site simple to have arbitrarily high type.  This uses
the bounded-type exclusion, not the conclusion of the forward theorem
being proved.

Now put \(E=K_J\).  The axis
\[
 g_D\vee X_{T'\setminus D_*}\le J
\]
supplies independent new-site supports in \(E\), modulo \(Z\).
Inside one original site take the product of all individual
simple-factor intersections of \((E\cap Y)/Z\).  This product is
\(E\)-normal.  It is missing relative to \(K_{g_*}/Z\): the axis
intersection just obtained is a product contained in it, and
supplementation recovers that axis label by multiplying its lift by
\(K_{g_D}\).  That axis is not below \(g_*\).

By \Cref{fa:pure-links} on \([K_{g_*},K_J]\), the supported normal
product has image containing the nonabelian monolith modulo the base
core.  It is nonsoluble, and transitivity makes every supported factor
nonsoluble.  If \(J<r_D\), each factor is proper: a full one would,
by the transitive \(K_{g_D}\)-action and supplementation, give
\(E=K_{r_D}\).  This proves precisely the supported-product conclusion
needed below; it does not identify that product with all of
\((E\cap Y)/Z\).  In direct-system applications the block centers
are retained, so lifting these supports with their centers gives
the required actual block-supported matrices.

\end{proof}

\begin{remark}\label{fa:quotient-labels}

When reasoning in the almost simple quotient at \(A\) we generally keep ordinary subgroup labels
(including \(H\)) for their images as well. Overgroups there containing the base correspond without
change of order to ordinary labels below \(K_{M_A}\), since the kernel was already in the base. The image
of \(B_A\) still identifies with the \emph{actual} \(B_A\) by injectivity. Thus even if normal-product
transport is used while reasoning in this quotient, the inner groups transported can be taken back inside
that actual simple group first. Ordinary labels throughout are products at the original sites; their
being a sublattice, rather than necessarily the full displayed label lattice, matters for cover/size
arguments. Within the ideal every relative cut \(C_D\) is product (global cut intersected with \(M_A\));
in particular a maximal ordinary proper label has to be a cut, by containment.

\end{remark}

\begin{lemma}[When a product cover is a full cover]\label{fa:product-cover}
Let \(\Delta\) act on a finite direct product by coordinate-permuting
automorphisms.  Let \(X=\prod_i X_i<Y=\prod_i Y_i\) be invariant
coordinate products, with no invariant coordinate product strictly
between them.  If an invariant subgroup \(U\) satisfies \(X<U<Y\),
then \(X_i\trianglelefteq Y_i\) at every site.  Consequently,
nonnormality at one site makes \([X,Y]^\Delta\) a cover in the full
invariant-subgroup interval as well.
\end{lemma}

\begin{proof}
The coordinate closure of \(U\) is \(Y\).  Conjugating the independent
copies of \(X_i\) inside \(U\), whose projections are all of \(Y_i\),
therefore gives
\[
 N=\prod_i\langle X_i^{Y_i}\rangle\le U.
\]
This is an invariant coordinate product with \(X\le N<Y\), so the
product-cover hypothesis forces \(N=X\).  Thus every \(X_i\) is
normal in \(Y_i\).
\end{proof}

We apply this lemma before taking ancestor-centralizer quotients.
An ordinary cover is a cover among the original-site product labels.
If a formal intermediate exists, the lemma forces normality of the
lower site projection in the upper one.  This normality persists on
passing to an ancestor quotient and intersecting with an actual
supported subgroup there.  The concrete nonnormality tests below can
therefore detect a formal intermediate even when that intermediate
is not itself a product label.

\begin{lemma}[The three cuts above a disjoint meet]\label{fa:cut-meet}
If \(D,D'\) are nonempty and disjoint, the only coatoms above \(C_D\wedge C_{D'}\) in the relative marker lattice are
\[
 C_D,\qquad C_{D'},\qquad C_{D\cup D'}.
\]
Each of the first two covers the meet, whereas the interval up to the third is the Young subgroup
interval \([S_D\times S_{D'},S_{D\cup D'}]\).  These three cuts and the first two cover relations persist
in the ordinary product sublattice.
\end{lemma}

\begin{proof}

In particular, in the relative label lattice, the meet of cuts on \emph{disjoint nonempty} \(D,D'\)
absorbs everything outside \(D\cup D'\), with side group \(S_D\times S_{D'}\). The only nonempty sets
available for coatom cuts above that meet are \(D,D',D\cup D'\) (unions of the unabsorbed orbits). The
first two coatoms cover the meet, since inside one such cut one would have to absorb from the other set,
hence absorb that whole set, for any strict increase. Up to the union cut instead we have the full Young
subgroup interval in \(S_{D\cup D'}\). These are formal statements before the product selection (and
within the ideal at \(M_A\)); in particular all three coatoms persist, and formal covers between products
are still covers. In the other direction when needing an intermediate such as the formal wreath label,
which need not itself be product, we will argue using product closure rather than just assuming
persistence.

\end{proof}

\begin{lemma}[Finite selection simultaneously at all ancestors]\label{fa:all-ancestors}
For the successive exclusions below, fix a bound on path-prefix size and a finite sequence of rounds, each testing boundedly many
omitted positions and finitely many structure or comparison predicates.  A sufficiently large marker
reservoir can be thinned so that, conditionally on absence of the point configuration, each round applies
simultaneously at every eligible prefix carrying the simple group of \Cref{fa:relative-labels}.  At any
such fixed prefix, a bad witness with omitted support of size \(k\) has replicas on every \(k\)-subset of
a later reservoir.  Any prescribed bounded unions of these replicas retain all properties established in
preceding rounds.  The required reservoir sizes depend only on these finite requirements, not on group
dimensions or field sizes.
\end{lemma}

\begin{proof}

Apply finite Ramsey homogenization to the bounded omission patterns. Concretely color bounded-size
ordered index subsets by all the tested label predicates (Boolean, also subdivided into the boundedly
many indicated cases) relative to their prefixes \(A\) whenever those have the required path-node normal
simple in the ordinary product interval. It is canonical if it exists, by its soluble centralizer.
Include type predicates and the bounded comparisons across subsets as used for selection. From a
sufficiently large homogeneous set, for fixed \(A\) with fresh indices after it, any bad label with small
nonempty omitted set after \(A\) thus gives replicas with \emph{every} \(k\)-set of later fresh positions
(\(k=|D|\)), including as many disjoint ones as needed. Choices of witness structures can also be further
finitely colored and thinned. The exclusions are used successively: earlier ones for larger bounded sets
as needed to make the intersections (which contain the axes omitting the support unions) common
previous-good labels, just as properties there passed upwards from anchors. Thus they apply
simultaneously relative to all indicated high-type ancestors. All bounds for combined supports in each
round need only be finite in terms of subsequent round requirements, and never involve the ambient group
dimensions or fields.

More explicitly one may use successively shorter \emph{initial segments} of a homogeneous reservoir as
rounds proceed. Assert each exclusion (conditional on absence of the point configuration) for all prefixes \(A\) of the
allowed path sizes with the stated high simple data, and omitted sets in that round's segment, \(A\)
before the further omissions in index order. If a bad full stabilizer contains a tested axis label, its
exact omitted support is a nonempty subset of the tested further omissions. Replicate this support at
fixed \(A\) using the spared tail of the preceding segment; earlier-good properties are then available
for each specified bounded union of the replicas by the result of preceding rounds on that segment. Thus
choose earlier omission bounds and segment lengths large enough \emph{backwards} (e.g. include combined
supports of the finitely many simultaneous intersections needed in a later round, and allow extra finite
Ramsey thinning there). Predicate colors on ordered subsets can include all splits/sublist choices at the
finitely bounded sizes, recording actual path-simple existence even before descending (\(B_A\) is unique
by its soluble centralizer), needed same-characteristic/type tests, and existence of a witness with
exactly those support positions in the relevant categories. Not the unbounded field sizes, group orders,
etc., only finitely many case/comparison predicates are needed. In a contradiction at a fixed prefix one
can further choose witnesses and color them by finitely many tests, without losing earlier-good
properties. All unions and links still refer to the original \(T\); in particular the residual marker
counts for the normal/missing-site tests remain arbitrarily large.

\end{proof}

\subsection{Alternating ancestors}
\label{fa:alternating}
Throughout this subsection \(B_A\) is alternating of sufficiently large degree.  The homogeneous
families, disjoint replicas, and inherited properties are supplied by \Cref{fa:all-ancestors}.  A
previous-good intersection means one to which the conclusions of all preceding pruning rounds apply.

\begin{proposition}[Subset stabilizers produce points]\label{fa:subset-points}
If a tested descendant is contained in a proper subset stabilizer, then, after the finite selections of
\Cref{fa:all-ancestors}, there are arbitrarily many distinct relative labels \(M_{\{i\}}\) that are full
stabilizers of individual natural points.  In absolute notation these are \(M_{A\cup\{i\}}\).
Consequently, in the absence of the conclusion of \Cref{r:forward-marker}, the required descendants and
their bounded intersections are transitive.
\end{proposition}

\begin{proof}

For subset stabilizers \(K_{r_D}\), among the two supported alternating groups precisely one side is
missing relative to \(K_{g_D}\) by the normal property of \Cref{fa:pure-links} (soluble quotient over
both). Use the retained side's point subset \(U_D\); \(K_{g_D}\) acts transitively there by
supplementation with the other side, even for small size. Thus retained point sets for disjoint supports
are disjoint, by mutual containment of each smaller label in the other's stabilizer and distinctness. In
a homogeneous \(k\)-set family we can take their sizes each below a small fixed fraction of the whole
(color, using arbitrarily many disjoint replicas). Their stabilizers are maximal: an increase would give
transitivity and overlapping conjugates of the large complementary supported alternating. So \(r_D=C_D\)
(relative labels; all relative coatoms are products by intersecting global cuts with \(M_A\)). The same
maximality applies to the subset union for two disjoint \(k\)-sets. The union label contains \(C_D\wedge
C_{D'}\), and therefore equals \(C_{D\cup D'}\), since the only coatoms above that meet are these three,
and the union stabilizer moves the individual subsets (correct parity outside). Consequently the union
itself is independent of the split by its small size and the stabilizer orbits.

\smallskip\noindent\emph{Decomposing the retained subsets.}

It follows that \(U_D\) splits into elementary disjoint subsets indexed by the individual \(i\in D\), the
same in every use. Indeed for any point the indicator \(p_D\) has \(p_D+p_{D'}\) fixed by the union of
disjoint index sets. Thus \(p_{F\cup i}-p_{F\cup j}\) on replacing one index is independent of the other
\(k-1\) indices \(F\), by comparing with a common disjoint auxiliary \((k-1)\)-set. For \(k\ge2\) we
obtain a sum of individual real weights plus a constant: the replacement differences obey the triangle
identity by using the same \(F\) for three positions, so are differences of weights, and replacements
connect all the \(k\)-sets. Weights differ only by zero or one up to sign; with our ample reservoir, two
or more each of a high and a low weight would produce a difference of two between indicator sums. So
either all indicators vanish (a constant one violates disjointness) or just one position is exceptional
and membership is exactly when it occurs (the reverse choice again violates disjointness). For \(k=1\)
use disjointness directly. At most one elementary subset can be empty by label distinctness. For any
nonempty one its stabilizer label is \(M_{\{i\}}\), since it contains \(C_D\wedge C_{D'}\) for pairs
meeting just in \(i\); two such meets with remaining positions disjoint already join to \(M_{\{i\}}\).
Indeed each meet fixes position \(i\) and absorbs all outside \(D\cup D'\); use four sets all meeting
pairwise just at \(i\). For \(k=1\) we already have the label \(C_{\{i\}}=M_{\{i\}}\) itself.

\smallskip\noindent\emph{Equal sizes and singletons.}

Unions of two or four of these nonempty elementary sets again have the corresponding cut labels by the
coatom-meet argument. For disjoint pairs \(P_1,P_2\), the full interval from the two cuts' meet \(J_0\)
to the union cut \(F_0\) is that from \(S_2\times S_2\) to \(S_4\), with a proper wreath intermediate. If
actual union sizes for the two differed, the ordinary interval would be a cover (project the full
stabilizers to the symmetric group on the combined set, using outside parity correction).
Apply \Cref{fa:product-cover} to the product endpoints in the original
power.  The formal wreath intermediate would force
\(K_{J_0}\cap S\trianglelefteq K_{F_0}\cap S\) before the ancestor
quotient.  Passing to that quotient and intersecting with the actual
\(B_A\) gives a contradiction: on the combined set the upper group
induces the full symmetric group by outside parity correction, whereas
the lower group induces exactly its two-part stabilizer.  For
unequal positive part sizes that stabilizer is maximal proper in the full symmetric group (an increase is
transitive, merging overlapping larger parts via conjugates of their supported symmetric group), not
normal. Thus pair-sums of sizes agree and all nonempty elementary sizes agree (compare pairs sharing one
index against a common disjoint pair). For two elementary sets the analogous formal interval is just a
cover, so their equal size must be one by block-swapping with parity correction. This gives the desired
point labels. Henceforth we can use transitivity after this exclusion if success was not already
obtained.

\end{proof}

\begin{lemma}[Exclusion of imprimitivity]\label{fa:imprimitive-exclusion}
Assume the transitivity conclusion of \Cref{fa:subset-points} for all bounded unions required in this
round.  A homogeneous family of proper full stabilizers of nontrivial equal-block partitions cannot occur
at the tested descendants.
\end{lemma}

\begin{proof}

For full stabilizers of nontrivial equal-block partitions (the imprimitivity test) let the sizes be \(a\)
per block, \(b'\) blocks. If \(a\ge3\) with \(A_a^{\,b'}\) retained, the supported primitive groups force
refinement into blocks of another system of the pattern (discrete restriction to a block is impossible by
moving a point while fixing its equivalent outside neighbor), hence equality mutually. For \(a=2\) use
instead retained even simultaneous flips of two blocks, precluding any other matching of this kind.
Otherwise \(a\ge5\) and the alternating base product is missing. Here the full stabilizer is maximal by
the supported \(A_a\)'s: an increase gives overlapping conjugate blocks, thus full alternating on merged
components, and its primitive top on blocks forces total merging. Indeed supported alternating groups of
these sizes on overlapping sets generate alternating on the union (conjugate three-cycles through a
common point; even with just one common point this brings in three-cycles using one point of the other
set and two of one's own). The full top symmetric action is available, with internal parity correction.
Take two disjoint failures, giving coatoms \(C_D,C_{D'}\) and a meet covered in each. The two block
systems have common cell size \(h\) on occupied incidences by prior transitivity. They cannot coarsen one
another by primitive top action, which is realized inside the intersection for both by supplementation
relative to the \(g\) labels. If \(h\ge2\), the cell-refinement stabilizer intersects a row label
properly (by its supported primitive block groups), strictly above the intersection because one can swap
two cells within a row while fixing points elsewhere, with internal parity correction, breaking the
columns. This contradicts covering. If \(h=1\), the intersection embeds subdirectly in \(S_{b'}\times
S_{b''}\) via the two top actions, both block counts at least five by the touching block sizes. The
kernel of projection to the first top action must still supply all even-total base sign patterns there
(supplementation with \(A_a^{\,b'}\), which acts trivially on that action's blocks and signs). But this
kernel injects as a normal subgroup of \(S_{b''}\); its abelianization has size at most two, versus its
required homomorphic sign image of size at least \(2^{b'-1}\).

\end{proof}

\begin{proposition}[Exclusion of affine and Cartesian structures]\label{fa:cartesian-exclusion}
After \Cref{fa:imprimitive-exclusion}, the tested descendants are primitive.  They cannot lie in full
normalizers of elementary abelian regular groups or in proper Cartesian-system stabilizers.
\end{proposition}

\begin{proof}

Using then common primitive previous-good labels \(E\), full normalizers of elementary abelian regular
groups are excluded by \Cref{c:alternating-affine}: the regular groups are retained and hence mutually
included in \(E\), where each is minimal normal; distinct such groups commute, whereas an abelian regular
group is self-centralizing.  For Cartesian stabilizers, \Cref{c:cartesian-retained} applies with mutual
inclusion in intersections for disjoint supports; small slot sizes \(2,3,4\) are likewise precluded by
the regular-group test. In the missing power case for two disjoint supports, the supported normal
products in \(E\) are supplied by \Cref{fa:intersection-support}, so the tuple commutator and
point-stabilizer factorization in \Cref{c:cartesian-refinement} give a common Cartesian refinement with
at most one nontrivial fine slot per pair of slot
positions. Take \(E\) to be the two-label intersection. Its primitivity gives transitivity on occupied
fine slots, say each of set size \(h\ge2\). The top action at each of the two original systems inside
\(E\) is transitive containing the alternating group by supplementation, hence primitive, so grouping
fine slots into rows and columns splits each slot position on both sides into at least two fine ones
(refinement would be improper by that primitivity). Both original slot set sizes are large by the
bounded-type consequence for the missing powers.

Let \(u\) be the intersection of the full row-system stabilizer with the fine-slot stabilizer. This is a
proper decrease by the supported natural alternating action on each row slot, and it contains \(E\). It
permits interchanging two fine slots in one row, leaving the others in place: if the parity of the
interchange needs correction then \(h\) must be odd (at least three fine slots), so use an internal
symbol transposition. Together with the full column-system stabilizer this generates the ambient group.
Here are details of the latter assertion. Local alternating gates on two overlapping sets of Cartesian
directions each of product size at least ten generate the alternating gates on their union. Split into
differences and shared part of alphabet sizes \(a,c\) and \(b\), all at least two (proper overlap). On
the union the generated group is two-transitive: ordered distinct pairs can be changed to any one fixed
pair distinct in all three positions using two-transitivity on each of the two touching position-sets
(first obtain a difference in the shared part on a position-set with some difference, then set the first
and shared values as desired, then the shared and last). For \(b\ge3\) use three-cycle gates on the two
position-sets whose lifted supports intersect in one tuple there (one shared value appearing once in each
cycle and disjoint other shared-part values). Their commutator supplies a three-cycle, sufficient by
primitivity. For \(b=2\), both \(a,c\ge5\). Flip the shared bit on two values of the first alphabet, and
use arbitrary \(A_c\) conditional on bit one. Commuting these and commuting again with conditional
\(A_c\) gives \(A_c\) conditional on bit one and membership in the chosen two values (by perfection).
Indeed the first commutation applies an element of \(A_c\) on one bit value and its inverse on the other
just for those two first-alphabet values; intersect the condition with bit one by commutators again. Two
such two-value choices intersecting once and commuting their conditional actions likewise give supported
\(A_c\) on a single fiber, again sufficient (a three-cycle there, then Jordan). All gates obtained act
independently of unused coordinates, meaning the same union permutation repeated on all outside tuples.
Now we have the column gates and those on the conjugate columns by the interchange. Indeed an altered
column and its original overlap properly, and merging them extends the latter by a slot from a second
column, next allowing merger with that entire column by the same observation. Align actual permutations
of the columns preserving fine coordinates up to a permutation (uniform fine-slot alphabets and common
fine-slot count per column); even column permutations so aligned are in the socle as well. Thus the
merged directions connect all columns by primitivity (for two columns the one interchange already
connects). The gate observation yields the whole alternating group, hence full join (the base already
supplements it). This contradicts \(u\vee r_{D_2}<M_A\).

\end{proof}

\begin{lemma}[Exclusion of primitive diagonal structures]\label{fa:diagonal-exclusion}
After the preceding alternating rounds, primitive diagonal structures are excluded as well.  Hence the
required descendants of an alternating ancestor are primitive almost simple groups.  The transitivity
conclusions continue to hold relative to all such ancestors.
\end{lemma}

\begin{proof}

By \Cref{c:primitive-diagonal}, every remaining primitive group that is not almost simple normalizes a
transitive power \(Y=B^s\), \(s\ge2\), with a full possibly twisted diagonal point stabilizer and natural
degree \(|B|^{s-1}\).  This includes two opposed regular simple factors.

First consider retention.  A retained power with transitive factor action is a minimal normal subgroup of
the common primitive intersection \(E\), with trivial centralizer because its diagonal point stabilizer
is self-normalizing.  Distinct retained powers would therefore coincide.  If the factor action is
intransitive, each factor-orbit product is transitive; in the diagonal action a transitive subproduct
uses at least \(s-1\) factors.  Thus the only possibility is \(s=2\), the two opposed regular simple
factors.  A retained regular factor is minimal normal in \(E\), and its centralizer supplies its opposite
regular group.  Distinct minimal normals commute and consequently determine the same opposed pair.  These
are exactly the retained alternatives in \Cref{c:diagonal-pruning}, and exclude distinct disjoint-support
witnesses.

It remains to consider a missing power with transitive factor action.  Take disjoint labels and their
proper previous-good intersection \(E\).  By \Cref{fa:intersection-support}, \(E\) has a missing normal
product \(R=\prod_{a=1}^s R_a\) of singly supported proper groups, with nonsoluble factors.  This
assertion does not require \(R=E\cap Y\).  Choose a minimal \(E\)-normal subgroup \(P'=C^u\) inside
\(R\).  It is nonabelian by the affine exclusion, and transitive by primitivity of \(E\).  Every factor
of \(P'\) lies in one \(R_a\): the groups \([P',R_a]\le P'\cap R_a\) are normal in \(P'\), distinct slots
commute, and these commutators generate \(P'\), since \(P'\le R\) is perfect.  The assignment is
\(E\)-equivariant.  Supplementation relative to \(g_D\) makes \(E\) transitive on the \(s\) slots, so
there is a common positive number \(h\) of factors over each slot and \(u=sh\).

The primitive-group argument of \Cref{c:primitive-diagonal}, applied to \(E\), says that the point
stabilizer \(Q=P'_x\) is a full diagonal.  If \(h>1\), the product of the projections of \(Q\) to the
\(s\) assigned batches is a proper intermediate subgroup between \(Q\) and \(P'\), normalized by \(E_x\),
contradicting primitivity.  If \(h=1\), equality of point degrees gives
\[
 |C|^{s-1}=|B|^{s-1}.
\]
But a factor \(C\) embeds in a proper subgroup \(R_a<B\), an impossibility.  Thus the diagonal case is
excluded, and the required descendants are primitive almost simple.  These conclusions apply
simultaneously at the eligible ancestors by \Cref{fa:all-ancestors}.

\end{proof}

\subsection{Classical subspaces and projector stabilizers}
\label{fa:classical-initial}
We use the natural covering groups, forms, collineations, and anti-actions of \Cref{sec:classical}.  All
dimensions required to be large can be made so by \Cref{fa:bounded-types}.

\begin{lemma}[Exclusion of parabolic stabilizers]\label{fa:parabolic-exclusion}
For a high classical ancestor, the tested descendants cannot preserve a proper linear subspace in the
absence of anti-actions, a proper signed flag in their presence, or a proper singular subspace in the
isometry cases.  The same assertion holds for the bounded intersections needed in subsequent rounds.
\end{lemma}

\begin{proof}

Soluble unipotent radicals are retained, thus each lies in all other labels on a disjoint family. For
linear subspaces (no anti-actions), signed flags (linear with anti-actions), or singular subspaces in the
isometry types, the shear arguments of \Cref{c:linear-subspaces,c:isometry-subspaces} give
a chain (using the smaller space of each signed flag, the other determined in reverse by an anti-action
of the base), up to the bounded singular-space terminal exception. Pick two in a many-member disjoint
family there (below terminal dimensions with ambiguities in orthogonal type, e.g. below Witt index minus
one in split type D). Their full stabilizers and intersection meet the ancestor socle in the
corresponding parabolics of compatible flags. By the standard parabolic overgroup theorem for the natural
\(BN\)-pair in the simple Lie-type group (overgroups of a parabolic are parabolic, with the
standard-parabolic subset order; in the classical interpretation use the indicated subspace or singular
flags, away from the split type \(D\) terminal fork here), the labels are maximal and their meet has only
the two maximal increases to coatoms: overgroups intersect the socle in parabolics obtained by dropping
flag terms, determining the overgroup since it contains the supplementing \(H\). In the signed case terms
must be kept in duality-stable pairs (possibly a single middle term). This contradicts the three coatoms
above the meet of two coatoms with nonempty disjoint supports.

\end{proof}

\begin{lemma}[Nonsingular radical lines in characteristic two]\label{fa:radical-lines}
In the orthogonal case in even characteristic, a full stabilizer of a nonsingular radical line is
maximal.  Such line stabilizers cannot occur in a homogeneous family of the tested descendants after
\Cref{fa:parabolic-exclusion}.
\end{lemma}

\begin{proof}

Indeed write the hyperplane of orthogonality to a line representative \(u\) as \(u^\perp\); the
simple-group stabilizer contains products of two nonsingular symmetries with centers therein. An increase
gives another such hyperplane \(w^\perp\); there is a common nonsingular partner for bridging the
products (high dimension). Odd and even words of symmetries using these two hyperplanes then give by
conjugacy all centers, so even products giving the simple group are contained. Indeed arbitrary even
words from the two center sets together are available by inserting a common partner for mixed products;
once all centers are obtained as symmetries by words (necessarily of odd length by Dickson parity), their
pair products thus belong and give a nontrivial normal generating set projectively by simplicity. In
detail, for nonsingular \(v\) outside both one can choose nonsingular \(z\in u^\perp\) whose symmetry
takes \(v\) into \(w^\perp\), by choosing both polar pairings \(b(z,v),b(z,w)\) nonzero there, then
ensuring
\(Q(z)=b(z,v)b(z,w)/b(v,w)\) for the quadratic form \(Q\). One can adjust this value arbitrarily with a
large nondegenerate subspace orthogonal to the indicated vectors. Neither functional vanishes identically
since the lines are distinct and \(v\notin u^\perp\). This proves the generation assertion, also
sufficient in the ambient quotient since the base supplements its socle. For two disjoint failures the
pair span \(P\) must be nondegenerate, else its singular line would be preserved and has been excluded on
the combined support. The pair intersection then agrees with the full intersection of either line
stabilizer with the plane stabilizer (by covering). Indeed the latter cuts it properly, e.g. products of
nonsingular symmetries in \(u^\perp\) move the plane (choose one center pairing nontrivially with \(w\)
and outside the plane, correcting by one fixing the plane). But this full intersection on the \(u\) side
includes the symmetry at \(u\) corrected outside \(P\), failing to preserve \(w\). This gives a
contradiction.

\end{proof}

\begin{proposition}[Maximality of an unbalanced ordered splitting]\label{fa:projector-maximal}
Let \(V=W\oplus U\), with \(0<\dim U<\dim W\), be a nondegenerate splitting in an isometry type, or an
image--kernel splitting in linear type with anti-actions.  Let \(M\) be a group of the corresponding
classical actions containing the natural special group, and let \(K\) be the full ordered splitting
stabilizer, equivalently the idempotent stabilizer for the ring action.  Assume that \(K\) supplements
the special group in \(M\), that its collineations act semilinearly irreducibly on \(U\), and that the
supported special group on \(W\) is available, as is that on \(U\) when \(\dim U\) is large.  In the
linear case assume also that \(K\) contains an anti-action fixing the projector, hence carrying its
image--kernel data into the dual annihilator data.  For sufficiently large \(\dim W\), the subgroup \(K\)
is maximal in \(M\).
\end{proposition}

\begin{proof}

Let an overgroup properly contain \(K\).  We show that the transported
copies of the supported special group on \(W\) generate the whole
ambient special group.  First we find large summands close to \(W\)
whose vectors span \(V\) in boundedly many summands.  We then collapse
those sums and recover the global elementary generators.

Throughout, supported groups use their natural matrices, equal to the
identity off the indicated summand; their projective images are the
groups in \(M\).  Transporting a splitting transports these groups
with it.  In linear type an anti-action transports the special linear
support by inverse transpose, and exchanges the corresponding primal
and dual data.

\smallskip\noindent\emph{Near splittings and bounded sums.}

Any increase contains transported copies of the supported special on \(W\) at the largest summand of any
moved projector. This copy cannot normalize the original projector: by irreducibility on its support and
dimensions it would first force its image space to agree with \(W\), then its fixed complement to be
\(U\). The high classical special groups here are generated by matrices of bounded displacement rank
(ordinary elementary/root generators, also obtained by Witt reduction). For instance take all inner
conjugates of a noncentral elementary generator (in an isometry type one can use a special linear shear
on two totally singular basis directions extended dually and identically off the associated two Witt
pairs); they suffice by projective simplicity and perfection (generation modulo the special group's
center suffices by perfection). Thus a near copy of the splitting with moved projector is attained, with
new largest space \(W'\) differing from \(W\) by bounded codimension of intersection. Use the anti-action
if necessary in linear type to get movement of the \emph{image} (also to achieve kernel movement). Use these
and their \(K\)-transports as near splittings (transports by anti-actions of these bounded-rank
displacements remain near). Every vector is a sum of boundedly many from these large spaces. Indeed
projections span \(U\) by semilinear irreducibility. If \(U\) is bounded-dimensional this suffices.
Otherwise transport a nonzero such projection by the available classical actions to get bounded-sum
generation there. In the linear or symplectic case use transitivity, and for a norm constraint use high
Witt dimension: every vector \(v\) is a sum of boundedly many vectors of a given nonzero vector's type
(norm, including isotropic nonzero). For instance prescribe \(y\) and \(v-y\) both of the given norm by
first imposing the corresponding polar pairing with \(v\ne0\) (a trace in hermitian type), then adjusting
the norm in a large orthogonal subspace, avoiding zero vectors. Here we use the usual transitivity on
nonzero vectors of given norm also under the high special by Witt extension with determinant, spinor or
Dickson corrections: a large nondegenerate subspace orthogonal to the vector to be fixed supplies
arbitrary such coset data (in unitary type norm-one determinants; in orthogonal type
reflections/symmetries at nonsingular vectors, both norm-square classes available if needed). Likewise
after prescribing the pairing to get two summands of the given norm, orthogonal adjustment can be made
perpendicular also to a first solution of the pairing, using a large nondegenerate subspace representing
every norm many times. Use \(W\) itself for the unprojected correction. In the linear setting dual
generation holds also: under transport by a fixed anti-action at the original splitting, large-space
images for near projectors correspond by a semilinear primal-to-dual identification to annihilators of
the transported kernels. The family of these initial near splittings can be taken closed under that
transport in both directions (bounded defect on both images and kernels). Thus these annihilators span
the dual, with the same bounded-sum conclusion there; in isometry type use selfadjointness and primal
spanning. In particular the common kernel of the available near projectors to the large spaces is zero.

\smallskip\noindent\emph{Collapsing bounded sums.}

One can collapse such bounded sums to one summand from some conjugate splitting, still within bounded
codimension of \(W\). Write \(W_i\) for summand spaces used and \(p_i\) for their projectors, with
bounded pairwise defect, so intersections also contain large nondegenerate spaces when relevant. If
\(s=w_i+w_j,\ w_i\in W_i,w_j\in W_j\), has \(p_i(s)\ne0\), choose \(d\in W_i\cap W_j\) so \(p_i(w_j+d)\)
is nonzero of matching type/norm as \(p_i(s)\). Use largeness of the intersection to solve the norm
condition, with nonzero choices, and the special-group orbit transitivity
with corrections established above.  The complementary components agree:
\[
 (1-p_i)(w_j+d)=(1-p_i)w_j=(1-p_i)s.
\]
The supported group on \(W_i\) can therefore send \(w_j+d\) to
\(s\), mapping \(W_j\) to another near splitting.  For sums of three or more, if all pair-collapse
projections vanish, perturb two summands oppositely by an intersection vector to make another pair
collapse. For example perturb the first and second by a common nonzero vector with opposite signs; the
first projection for the first-plus-third sum is now nonzero. Defects stay boundedly large at most
because only boundedly many steps are needed.  More precisely, write
\(\delta(E,F)=\dim E-\dim(E\cap F)\) for these equal-dimensional
spaces.  A transport \(g\) supported on \(W_i\) fixes its projector,
so
\[
 \delta(gW_j,W_i)=\delta(W_j,W_i),\qquad
 \delta(gW_j,W)\le 3d_0
\]
when both current defects from \(W\) are at most \(d_0\).
The same calculation holds on the dual supported spaces in linear
type, since \(g\) fixes that part of its splitting as well.
Thus we maintain bounded defects for both images and kernels.
Large nondegenerate parts of intersections in
these steps follow since intersection with a near space has bounded codimension in the nondegenerate
summand (and we may still take boundedly many orthogonality constraints). Finally if a nonzero pair sum
lies in both projection kernels, take a fresh near \(W_3\) with
nonzero projection of the sum.  For a nonzero summand \(w_i\),
choose in \(W_i\cap W_3\) a nonzero vector of its type and transport
it to \(w_i\) by the supported special group on \(W_i\).
This transport fixes the sum, since \(p_i(s)=0\).  Its image of
the \(W_3\)-splitting therefore covers \(w_i\) and still has
nonzero projection of \(s\).  Use that splitting to collapse the
pair.  The same construction holds dually in linear type, with
both defect bounds maintained as above.

\smallskip\noindent\emph{Global elementary generators.}

Thus all centers are available near \(W\). Supported special transvections now generate globally in
unitary and symplectic type. In orthogonal type use products of reflections/symmetries of same norm,
bridging by a common partner in the large near-space intersections; such products generate (in particular
supply global nontrivial normal generators). In linear type for \(v,f\) covered respectively in some
\(W_i\) and some dual supported \(W_j^*\), \(f(v)=0\), take a commutator of supported transvections with
tensors \(v\otimes h,\ z\otimes f\), with \(f(z)=h(v)=0,\ h(z)\ne0\), \(h\) supported in the first dual
space and \(z\in W_j\). These choices exist by the large intersections relative to the first projection
kernel (one can take \(z\) in \(W_j\cap W_i\cap\ker f\) independent of \(v\), then choose \(h\) on
\(W_i\); its scaling gives arbitrary transvection coefficient): in the linear argument keep bounded
intersection defects simultaneously on images in primal and dual. This holds for the initial near
splittings by bounded displacement (also of the contragredients); when a supported group transports
another near splitting in a collapse, it keeps its own splitting fixed, so both defects still grow only
boundedly. This gives the required full transvections. This proves maximality.\qedhere

\end{proof}

\begin{lemma}[Orthogonality of the retained projectors]\label{fa:projector-orthogonal}
After \Cref{fa:parabolic-exclusion,fa:radical-lines}, consider full stabilizers of proper nondegenerate
subspaces in isometry type, or of proper nontrivial individual projectors in linear type with
anti-actions.  Let \(e_D\) be the projector onto the side whose supported special group is retained
relative to \(g_D\).  After finite selection, witnesses on disjoint omitted supports satisfy
\(e_De_{D'}=e_{D'}e_D=0\).  For the small orthogonal exceptions the nondegenerate-line round is performed
first.
\end{lemma}

\begin{proof}

Now take isometry stabilizers of proper nondegenerate subspaces, or linear-with-anti individual proper
nontrivial projector stabilizers. Designate by \(U_D\) the side with supported special retained by the
normal property of \Cref{fa:pure-links} over \(g_D\); exactly one side is missing. Write \(e_D\) for
the projector to \(U_D\). For a pattern with large retained dimensions, disjoint supports then have
mutually orthogonal projectors as in \Cref{c:linear-subspaces,c:isometry-subspaces}, by supported
irreducibility and mutual preservation. For a pattern with bounded retained dimension (take it constant),
the full stabilizers used are maximal by \Cref{fa:projector-maximal}. We explain small
exceptions/ordering in isometry type: the full individual isometry group on \(U_D\) is available in
projection with corrections on the large complement. It acts irreducibly except possibly split orthogonal
planes over two and three. Over two there is a forbidden nonsingular radical line. Over three there are
two nonsingular axes individually invariant unless interchanged by a similarity giving semilinear
irreducibility. Thus first run the nondegenerate line exclusion, then these other dimensions with
invariant-line cases forbidden on the combined supports. There is no dependence of the line test itself
on excluding these planes.

For two disjoint bounded-dimensional occurrences the sum of the small spaces is again complemented, by
its orthogonal space or by the kernel intersection, via the prior singular/signed flag exclusions (also
nonsingular radical lines when relevant). In linear type an anti-element of \(H\) fixing both projectors
exchanges the sum with the kernel intersection, as do signed anti-actions of the pair group; a nonzero
intersection gives a signed flag. Also images determine kernels in this case. The sum projector
stabilizer cuts each individual stabilizer properly (large complementary supported group) and thus its
full intersection with either equals the pair group by covering. Consequently all linear or isometric
maps on one small summand, identity on the rest of the sum (split it using that summand's prescribed
complement) and corrected outside, preserve the other small space. The outside here is in both projector
kernels, giving a large space for arbitrary special-group coset corrections as above. Distinct small
spaces (complements determined and dimensions taken constant) cannot contain each other; this also
explains strict cutting by the sum stabilizer, since a projection of the sum to an individual large
complement is then nonzero proper, not preserved by the full supported large action. Taking differences
on the other small space under the indicated maps thus forces zero projection onto the chosen summand
whenever its individual linear/isometry group acts irreducibly and nontrivially: a nonzero projection
generates the irreducible space already by its orbit differences. Hence mutual orthogonality outside the
indicated small exceptions. For the split orthogonal plane over three with interchange, the isometries
here together with the intersection label interchanging the nonsingular axes by earlier line exclusion
suffice (the minus sign supported there on the sum splits the other small space along the projection, and
the joint action is irreducible on the summand by the interchange together with its isometries). For a
binary linear line with anti-action, use the bilinear pairing giving a fixed anti-element: each line and
kernel are mutual annihilators for this pairing, with self-value 1 at the line representative. On two
distinct such representatives the two cross-values agree, since left and right orthogonality to each line
coincide; value 1 off-diagonal would violate the complementation of the sum. This likewise gives
projector orthogonality.

\end{proof}

The orthogonal projectors admit an additive decomposition over the individual marker positions.  We
record the algebraic step separately because it works in every characteristic.

\begin{lemma}[The projector fold]\label{fa:projector-fold}
Let \(e_D\) be a family of idempotents in one associative algebra,
indexed by the \(k\)-subsets of a sufficiently large finite reservoir.
Suppose that idempotents on disjoint sets annihilate one another on both sides, and that
\(e_D+e_{D'}\), for disjoint \(D,D'\), depends only on \(D\cup D'\).  Then there are pairwise
annihilating idempotents \(h_i\) such that
\[
 e_D=\sum_{i\in D}h_i.
\]
If the \(e_D\) are selfadjoint, so are the \(h_i\).
\end{lemma}

\begin{proof}

For \(k\ge2\), \(\Delta_{ij}=e_{F\cup i}-e_{F\cup j}\) is independent of \(F\) as in the subset argument.
Disjoint-pair differences annihilate each other. Also \(e_D\) acts as identity on \(\Delta_{ij}\) on both
sides when \(i,j\in D\), and as zero when \(i,j\notin D\). The latter follows using a disjoint witness
for the difference, the former by adjoining a spare disjoint \(k\)-set to \(D\) (the union projector
supports a witness within the union). Then \(h_i=\Delta_{ij}\Delta_{il}\) is independent of distinct
auxiliary \(j,l\), including order, by disjoint annihilation.
For example, with a fourth index \(a\), the triangle identities give
\[
 \Delta_{ij}\Delta_{il}
 =\Delta_{ij}\Delta_{ia}
 =\Delta_{il}\Delta_{ia}
 =\Delta_{il}\Delta_{ij}.
\]
It meets each \(e_D\) by multiplication
either as identity-action (\(i\in D\), use a partner in \(D\)) or zero. To obtain
\(\Delta_{ij}=h_i-h_j\), choose \(D\) containing \(i,j\) and let \(D'\) replace them by two fresh
positions \(a,b\).  Then
\[
 \Delta_{ij}=\Delta_{ij}(e_D-e_{D'})
 =\Delta_{ij}(\Delta_{ia}+\Delta_{jb})=h_i-h_j.
\]
No division by the characteristic or by \(k\) occurs. Thus distinct \(h_i\)'s annihilate each other (separate by an \(e_D\)), and
\(h_i=(h_i-h_j)(h_i-h_l)=h_i^2\). They are orthogonal idempotents, selfadjoint when relevant, and the
remainder in \(e_D-\sum_{i\in D}h_i\) is constant, necessarily zero by disjointness and idempotence (each
such remainder is itself an idempotent supported by \(e_D\) on both sides). For \(k=1\) use the
projectors directly.

\end{proof}

\begin{proposition}[Exclusion of proper projectors]\label{fa:projector-exclusion}
Under the preceding classical rounds, the tested descendants cannot lie in any of the proper stabilizers
considered in \Cref{fa:projector-orthogonal}.  The remaining descendants and the required bounded
intersections therefore have irreducible anchors; in linear type their collineation parts are also
irreducible.
\end{proposition}

\begin{proof}

In the resulting homogeneous \(k\)-set patterns choose dimensions small relative to total by
orthogonality on disjoint supports. The stabilizers and disjoint-pair sum stabilizers are maximal by
\Cref{fa:projector-maximal}, and the sum label is \(C_{D\cup D'}\) as in the subset argument (it moves
each summand). For clarity the aggregate here has the requisite irreducible collineation action: split
planes over two do not arise as sums, and over three in the line test choose both norms of the same class
by coloring, avoiding the split plane. Thus the aggregate projector \(e_D+e_{D'}\) depends only on the
union; it is recoverable from its stabilizer, e.g. the supported large complementary special forces
uniqueness among projectors with small image.

Apply \Cref{fa:projector-fold} to write the projectors as sums of the \(h_i\).
In this marker configuration, at most one \(h_i\) vanishes,
and each nonzero \(h_i\) has full stabilizer label \(M_{\{i\}}\).
For \(k=1\), the full stabilizers are already the singleton
cuts \(C_{\{i\}}=M_{\{i\}}\). For \(k\ge2\), two vanishing
\(h_i\)'s would identify distinct cut projectors by replacing one index
with the other. For a nonzero \(h_i\), choose \(D,E\) meeting just
at \(i\); then \(h_i=e_De_E\). Thus the paired cut-meet argument
of \Cref{fa:subset-points} puts \(M_{\{i\}}\) in its full stabilizer.
That stabilizer is proper because \(h_i\) is a nonzero proper
projector, with nonzero disjoint witnesses remaining. Since
\(M_{\{i\}}\) is a coatom, it is the full stabilizer.

For nonzero positions, in large fixed sizes \(l,2l\) with ample further spares, all sum-projector
stabilizers on those sets \(P\) are again maximal (keep dimensions small relative to the ambient by
thinning), hence are \(C_{B(P)}\) for some nonempty \(B(P)\subseteq P\), and distinct by projector
uniqueness. Homogeneity/counting forces \(B(P)=P\) on a thinned reservoir (homogenize the decreasing-size
test; injectivity rules out strict decrease everywhere for that size). For disjoint \(l\)-sets, if the
sum dimensions differ, their meet to the combined cut is an ordinary cover, by
\Cref{fa:projector-maximal} applied within the combined aggregate (both pieces now sufficiently large).
Indeed projection to ring actions on that aggregate identifies the meet as full ordered splitting
stabilizer, with anti-action still present in the signed linear case. These applications have the
supplement hypothesis: the ancestor outer actions are supplied by the common base, and restrictions of
linear/special isometries from the inner ancestor are covered modulo the aggregate special group already
by the splitting stabilizer, using the two large pieces for arbitrary determinant etc. (orthogonal
determinant and spinor or Dickson corrections by Witt groups) and the large outside space for
compensation.  The formal wreath label is strictly between the two-cut
meet and the combined cut.  If these product endpoints formed an
ordinary cover, \Cref{fa:product-cover}, applied before the ancestor
quotient, would force normality at each original site.  This is
impossible on intersecting with the actual supported high special
group of the combined aggregate: its projective action is simple,
and the lower intersection is a proper nontrivial splitting stabilizer.
The ancestor centralizer kernel is already in the lower endpoint,
so this is also a valid nonnormality test in the quotient.  Hence dimensions agree for
disjoint \(l\)-sets, thus all individual dimensions agree (compare a pair of replacements to one common
disjoint set). Align the individual summand forms themselves by finite coloring. Now on an \(l\)-set the
sum-cut interval from the axis has at least the whole subgroup lattice of \(S_l\) just up to the
stabilizer allowing permutations of the individual projectors (kernel the axis label, and align
permutations with outside correction). The sum stabilizer exceeds this by its large full classical
action, which cannot preserve that uniform direct decomposition (same orders as in the earlier
direct-system pruning). Indeed the formal interval \([M_P,C_P]\) is exactly the subgroup lattice of
\(S_P\), with \(|P|=l\).  Its ordinary product sublattice cannot contain more labels.  Yet the
permutation stabilizer just described already supplies one ordinary label for every subgroup of \(S_l\),
and the strictly larger full sum stabilizer supplies an additional label.  This is the required
contradiction.  Thus these failures too are excluded. As
before we get prior irreducible anchors (collineation part also in linear type).

\end{proof}

\subsection{Direct systems}
\label{fa:direct-systems}

Use the direct systems and supported groups of
\Cref{c:direct-lines,c:direct-block-groups,c:direct-compatibility}, now with disjoint label
supports \(D_i\), \(K_i=K_{r_{D_i}}\) (as needed also for the following categories). Retention in the
smaller \(g_{D_i}\) labels places a group inside \emph{every} label \(K_j\) here. A previous-good anchor can
be taken to be the intersection or a large axis group below it, omitting just the joint support (all
relative to the ancestor). Thus we have transitivity and semilinear irreducibility on blocks by the
collineation part of the intersection whenever using them.

\begin{lemma}[Coincidence of exceptional matchings]\label{fa:direct-matchings}
In the retained line-system and exceptional plane-system cases of classical pruning, a family of
disjoint-support witnesses forces all the exceptional matchings to coincide.  A fixed previous-good
intersection then allows only boundedly many further systems.  Thus these cases are excluded by taking
sufficiently many replicas.
\end{lemma}

\begin{proof}

In the line-system argument of \Cref{c:direct-lines}, the diagonal groups are still mutually included by
retention, so it remains to treat its matching cases. All matchings relative to the
first system would have to agree: the signs of one of the other systems already interchange first-system
lines on a matched pair, simultaneously on a second pair if needed (choose two flips of equal norm class
for determinant/spinor correction). There are plenty of possible second pairs by transitivity and at most
two norm types (the two sign-orbit lines on a pair are of the same norm class); thus another invariant
matching must pair these same two first-system lines (otherwise choose the extra flip away from a
proposed different partner to keep that partner fixed). The choices on matched lines are unique as
before. For the exceptional plane matchings of \Cref{c:direct-block-groups}, an individual retained plane
sign likewise interchanges the
two old planes, forcing matching coincidence. The common intersection of all the labels is transitive on
the matching edges, so for that fixed intersection group bounded choices on one edge bound the number of
further systems, impossible. Here and below common intersections for any bounded family can be required
previous-good by earlier rounds at sufficiently increased omission bounds.

\end{proof}

\begin{lemma}[Support and compatibility for direct systems]\label{fa:direct-compatible}
In the remaining direct-system cases, two disjoint-support witnesses have the supported block groups
needed for the displacement-rank compatibility argument.  The block groups can be chosen as actual
block-supported matrices, with retained centers included.  Retained systems are excluded by comparable
refinements and the uniform order bound; a missing simple-site orbit has large block dimension and
bounded pair-orbit count.
\end{lemma}

\begin{proof}

In the remaining direct-system cases the supported groups of order greater than two for compatibility are
available exactly as above: use the full retained groups, or, for a missing simple-site orbit product,
\Cref{fa:intersection-support} with the previous-good intersection. (If two site orbits, can
likewise use full supports from a retained one.) Centers included with these lifts are retained. In fact
any missing simple orbit here forces large block dimension by the bounded-type consequence, with the
action on sites of bounded pair rank by supplementation relative to \(g_{D_i}\); in particular bounded
rank including the possible split four-space ambiguity cannot then occur. Block pair-orbits are bounded
by the alternating actions on at most two isometry classes (one class in linear type), and the simple
sites have fibers of size at most two over blocks. In applying the missing-product argument here the
singly supported projections modulo centers in the intersection actually yield block-supported matrices
projectively there: use representatives on that block in the supported choice of groups, since the center
kernel needed when lifting already lies in the intersection. We can take all available such block
matrices there, so the intersection's semilinear collineation stabilizer of a block normalizes the
corresponding supported group. The block module for the supported group again is semisimple without
trivial constituents by semilinear irreducibility for that stabilizer (use the module socle and the fixed
group of vectors, as before) and nontriviality of the block action. (For a missing product the singly
supported subgroups are even nonsoluble.) The displacement-rank proof of \Cref{c:direct-compatibility}
now applies: fixation of all the other system's blocks forces splitting by projections using orbit
differences there; otherwise movement both ways forces equal block dimensions, and forces any moving
block matrix to interchange a single pair of those other blocks as an involution, identity elsewhere
there and with displacement of full rank on its own support, as in \Cref{c:direct-compatibility}.  Its
retained-case proof also applies: full retained actions force comparable refinements, whose height is
absolutely bounded within the pattern by the uniform order obstruction on coarse blocks.

\end{proof}

\begin{lemma}[Merging compatible large summands]\label{fa:direct-merging}
Let \(W_1,W_2\) be compatible summands obtained by grouping blocks of one common direct splitting,
orthogonal in the isometry cases, with the corresponding complements.  If \(W_1\cap W_2\) is sufficiently
large and nondegenerate when relevant, their two supported special groups generate the supported special
group on \(W_1+W_2\).
\end{lemma}

\begin{proof}

Here the summands are direct groups of blocks of a common splitting, orthogonal in isometry cases, with
the corresponding complements. Use \(W_1\) of dimension at least \(\dim W_2\). Every vector in the sum
lies in some conjugate of \(W_1\) by the supported group of \(W_2\): for a nonzero \(W_2\)-projection
attain it by transporting a vector of the intersection (matching norm using its high rank), and the
remaining component already lies in \(W_1\). Thus we get the full symplectic or unitary transvection
generation on the sum; likewise products of orthogonal symmetries/reflections of matching norm, using
\(W_2\) and partners in its large intersections with these conjugates to bridge any two centers. These
elements give noncentral normal generation of the special. For linear type we also have dual coverage; as
in the projector argument get a transvection tensor \(v\otimes f,\ f(v)=0\), by commutators from two
supported spaces of dimension \(\dim W_1\), sufficient because twice that dimension exceeds the sum
dimension by a large amount. All assertions can use the natural covering matrices (the elementary
generators obtained give the perfect special group), with the indicated support.

\end{proof}

\begin{proposition}[Exclusion of the remaining direct systems]\label{fa:direct-exclusion}
After the preceding classical rounds, a homogeneous family of proper full direct-system stabilizers
cannot occur at the tested descendants.
\end{proposition}

\begin{proof}

Treat missing systems with two blocks first. More generally once these are excluded, two-class block
systems are impossible by coarsening, and the top action on blocks, induced inside intersections by
supplementation, is primitive (two blocks, or at least the alternating permutations with transitivity).
Thus two disjoint-support instances give incomparable compatible systems and each block splits into at
least two cells, of uniform dimension by transitivity. In the general round all those cells must in turn
be mutually isometric when appropriate (else a two-class coarsening). In the two-block round there is a
full \(2\times2\) grid of spaces, and for some choice of the pair of instances we can interchange two
same-row or same-column cells isometrically. Indeed if this fails the form must be orthogonal in odd
characteristic with odd cell dimension by similarity/transitivity. Take three failures, choosing their
common intersection to be previous-good as allowed by \Cref{fa:all-ancestors}.  Its collineation
irreducibility makes the occupied triple cells one orbit.  Since its permutation action lies in
\(S_2^3\), this orbit has size \(4\) or \(8\), and it projects onto all four cells for every pair. Size 8
violates odd pair-cell dimension by summing two equal triple-cell
dimensions. For size \(4\), the third two-block system is the unordered diagonal partition of the four cells of the first pair
(there is one third position at each occupied pair, with both third positions represented in each row and
column), so the whole intersection of that pair of stabilizers preserves it, impossible by disjoint label
supports.

Use a resolved pair, designated row and column, and let \(u\) be the intersection of the full row-system
stabilizer with the stabilizer of the simultaneous refinement. It contains the joint intersection and is
strictly smaller than the row label (uniform finer decomposition on each high-dimensional block). It
permits a swap of two fine cells in a row, leaving the others in place, or a row three-cycle of cells.
Indeed we have isometric choices as just described, and determinant, spinor or Dickson corrections inside
the cells when of large dimension. For small cell dimension here the row block dimension is large
(absolute thresholds can be made large enough to ensure at least three cells) and all cells are
equivalent, so take a three-cycle via commutators of aligned permutations. Such collineations come from
the ancestor socle.

We now apply \Cref{fa:direct-merging} to the original and altered column spaces.

A column space and its altered version lose/gain one cell by the within-row permutation and have large
intersection (at least half a block); their supported groups thus merge, extending the column \(X\) by a
cell from another column \(Y\). If that cell is large we merge with \(Y\) immediately. Otherwise use the
column label to transport within \(Y\) (alternating fine-cell permutations there while preserving \(X\))
and merge extensions along \(X\) until acquiring \(X+Y\). The resulting merged column pairs propagate to
a connected graph by block primitivity, and give full generation by repeated mergers with at least a
column in common. This contradicts the formal disjoint-support rule \(u\vee r_{D_{\rm col}}<M_A\), and
finishes direct systems.

\end{proof}

\subsection{Large tensor factors and representation bounds}
\label{fa:overgroups}

We supply maximality details for the forward field and ordered-factor tests, to avoid assumptions about
arbitrary structure stabilizers being maximal. Use ambient notation \(S\le M\le{\rm Aut}(S)\) in the
almost simple quotient at an ancestor, natural space \(V\) of dimension \(v\), \(\mathcal A={\rm
End}_F(V)\). We only use full stabilizers in groups supplementing \(S\) already by the common base. Write
\(Y\) below for the indicated supported projective special group(s), using their perfect matrix lifts
(center need not disappear in the projective field and paired constructions).

\begin{lemma}[Large-rank projective degree bounds]\label{fa:degree-gap}
For a sufficiently high-rank classical simple group of natural dimension \(m\), with the type and
characteristic conventions of \Cref{w:classical-conventions}, a nontrivial irreducible
projective representation in defining characteristic has degree at least \(m\).  Unless it is natural up
to Frobenius and duality, its degree is at least
\[
 \frac{m^2}{2}-O(m).
\]
A highest weight with \(l\) nonzero digit places has degree at least \(m^l\).  For field parameter \(P\),
every nontrivial cross-characteristic projective representation has degree at least \(P^{cm}\), for an
absolute \(c>0\); the same exponential lower bound holds for permutation degree.  These lower bounds on
nontrivial contributions apply to perfect central lifts and quotients.  Finally, a high-degree
alternating simple group acting projectively nontrivially in dimension \(v\) has natural alternating
degree \(O(v)\).
\end{lemma}

\begin{proof}

In defining characteristic, for classical simple groups of sufficiently high rank, with natural dimension
\(m\), the minimum nontrivial irreducible projective degree is at least \(m\); except for the natural
module up to Frobenius/duality, it is at least \(m^2/2-O(m)\). To see this use the restricted weights on
the universal cover and the Steinberg theory recalled in \Cref{w:classical-conventions}, using
\citet{Steinberg1963,Jantzen2003}. For a restricted digit in A the permutation orbit already
gives the latter bound except at a single endpoint only; at that endpoint with coefficient greater than
one, simple-root lowering (coefficient less than \(p\)) and the resulting orbit give it as well. In B,C,D
use signed coordinate permutations (even signs in D); at least two nonzero coordinates give the larger
bound in high rank (including spin weights), and a single first-coordinate support with coefficient
greater than one works again by lowering. Use B only in odd characteristic. Here the single first-node
coefficient one is the irreducible natural module, of degree \(m\) also in odd-dimensional B. For
\(k\ge2\) nonzero coordinates out of \(e\) in a signed word the orbit has size at least \(2^k\binom e
k\), allowing a loss of a factor two just at \(k=e\) with even-sign restriction; the stated quadratic
bound follows (including at \(k=2\)). These arguments also give the bound by \(m^l\) for \(l\) nonzero
digit places. Cross-characteristic nontrivial projective degree is at least \(P^{c m}\) for field
parameter \(P\) and absolute \(c>0\), by the argument of \Cref{w:characteristic-switch}; in particular the same exponential bound
works for permutation degree. These bounds on nontrivial contributions work also upon restriction to
perfect central lifts/quotients, pulling to the universal covers: lift projective actions there and take
a nontrivial composition factor (there is one by perfection if the projective action is nontrivial). The
center acts scalarly on irreducibles. In particular minimum permutation degree can be tested this way by
a permutation representation in another characteristic (the high perfect central lifts themselves then
have no nontrivial small permutation image). We use the classification and covering-group structure in
\Cref{w:classical-conventions}. An alternating simple of high degree acting projectively nontrivially in
degree \(v\)
has alternating degree \(O(v)\), by the following elementary argument.
Write its degree as \(n\), choose \(\ell\in\{3,5\}\) different from the
representation characteristic, and take \(r=\lfloor(n-2)/\ell\rfloor\)
disjoint \(\ell\)-cycles.  They generate \(E\cong C_\ell^r\), leaving
two spare letters.  The normalizer in \(A_n\) induces every permutation
of these generators: a permutation of the \(\ell\)-letter blocks has
its parity corrected, when necessary, on the two spare letters.
For matrix lifts of the projective action of \(E\), scalar commutators
define an alternating bicharacter \(b:E\times E\to\mu_\ell\).
Conjugation preserves \(b\).  Interchanging any two generators therefore
makes their value equal its inverse; because \(\ell\) is odd, that
value is one.  The lifts commute.  Rescaling them makes their
\(\ell\)-th powers one, and, since \(\ell\) is prime to the
characteristic, they are simultaneously diagonalizable.
The projective action of \(A_n\) is faithful by simplicity, so its
restriction embeds \(E\) in the \(\ell\)-torsion of the diagonal torus
modulo scalars.  Consequently \(r\le v-1\), and \(n\le5v+1\).

\end{proof}

\begin{lemma}[Coefficient twists on finite covering groups]\label{fa:coefficient-twists}
Let a classical finite covering group be defined by a standard pinned graph--field map of height \(f_0\).
Coefficient twisting a defining-characteristic irreducible representation by \(p^j\) is equivalent, after
restriction, to twisting its rational highest-weight lift by the source split \(p^j\)-Frobenius.  Thus
the nonzero digit positions are shifted by \(j\) modulo \(f_0\), with the diagram action at each wrap.
The same comparison may be made on composition factors.  Equivalence of fixed projective actions on a
perfect covering group implies linear equivalence of their lifts.
\end{lemma}

\begin{proof}

With a standard pinned graph/field map of height \(f_0\), coefficient \(p^j\)-twisting a
defining-characteristic irreducible is equivalent on the finite cover to twisting its rational
highest-weight lift by the split \(p^j\)-Frobenius of the source before restriction. Indeed in split
pinned coordinates, transforming coefficients in the rational representation together with the
coefficients in the argument (i.e. conjugating by coefficient Frobenius on points of target and source)
still gives an equivalent rational irreducible by highest weight. The split map permutes the standard
pinned finite fixed group; as in \Cref{w:classical-conventions}, digit positions thus move by \(j\) modulo \(f_0\), with
diagram action upon wrapping. This reasoning up to equivalence suffices (also for testing composition
factors); on a perfect covering group an equivalence of fixed projective actions lifts to linear
equivalence because scalar discrepancies are characters.

\end{proof}

\begin{proposition}[Maximality of an unequal ordered tensor factor]\label{fa:ordered-factor-maximal}
Let \(S\) be a high-rank classical simple group on \(V\), and let an ordered proper matrix tensor factor
give natural dimensions \(x<y\), both sufficiently large.  In isometry type assume that the factor is
selfadjoint.  In characteristic two assume that the bilateral factor forms are alternating, and in
orthogonal type use the canonical tensor quadratic form.  Exclude symplectic type in characteristic two.
Then the full stabilizer in \(S\) of the ordered factor is maximal.
\end{proposition}

\begin{proof}

Let \(Y_1,Y_2\) be the supported perfect projective special groups on the
two spaces, with the alternating special groups (symplectic groups) as needed on factors. They lie in
\(S\). For an overgroup within \(S\) take a minimal nontrivial normal subgroup \(N\) there. Its lift
algebra is semisimple by absolute irreducibility of the supporting action \(Y_1Y_2\), and has no multiple
center components: the larger simple action fixes all projectors by its minimum permutation degree, then
there are at most \(x\) of them (invariant spaces are modules over the larger matrix factor), so the
smaller also fixes them. The center is then scalar by perfection and full spanning by \(Y_1,Y_2\).

Any proper nonscalar matrix factor fixed individually by \(Y_1,Y_2\) gives precisely the original
factorization. Indeed lift the projective actions on factor spaces and tensor by perfection; all are
irreducible on the universal covers jointly by tensor irreducibility. The lower bounds \(x,y\) for any
respective nontrivial contributions (and \(xy\) if both in one slot) force their allocation to distinct
unique slots of a nontrivial split. Their matrix spans recover the factors. Here the separate lifts
commute as needed (any residual scalar commutators against a perfect factor give a trivial character);
joint irreducibles of the direct product over the closure are outer tensors. The supported matrices
originally span the actual factors over \(F\) by natural absolute irreducibility, and a trivial
projective action on other tensor directions means these matrices belong to the indicated one by
commuting with the others. Even for a permuted tensor system, \(Y_2\) on the larger space fixes all
slots, and acts projectively on just one (two would cost at least \(y^2\)); \(Y_1\) preserves that one
and must fix the others since their product dimension is at most \(x\). Thus a proper nonscalar lift
algebra for \(N\) already forces ordered stabilization by the overgroup. If the algebra is full and \(N\)
is a power of nonabelian simples, their lift algebras tensor to it by commuting as usual, so multiple
components also force ordered stabilization (no swap of unequal degrees). The elementary abelian
full-algebra case has the nondegenerate scalar commutator pairing, order \(v^2\) and projective linear
centralizer just itself. Its normalizer inside the simple group is far too small for \(Y_1Y_2\) (at most
\(v^{O(\log v)}\) by automorphisms of the elementary abelian group, while \(v\le y^2\)).

Thus otherwise \(N=B\) is a single simple with trivial projective linear centralizer and contains
\(Y_1,Y_2\) by Schreier. It must be high classical in the same characteristic: bounded Lie ranks are
excluded by the minimum faithful projective/linear degree of the larger simple support in any
characteristic (use also bounded-dimension adjoint realizations there); alternating type and
cross-characteristic classical type are excluded by the exponential bounds on embedding the larger
support, together with the action of \(B\) in degree \(xy\). In the alternating alternative the
alternating degree would be at most \(O(xy)\) by that action but at least exponential in \(y\) by
restriction to \(Y_2\). In the other-characteristic high classical alternative the natural dimension
there would already be at least exponential in \(y\) and its own cross-characteristic degree far exceeds
\(xy\). Its natural dimension \(w\) is at least \(x+y\): restrict its projective natural action to the
commuting supports, lift from universal covers of each (commuting linearly by perfection), and count
composition factors with nontrivial contribution from each (a joint irreducible with both nontrivial is
an external tensor over the closure, costing at least \(xy\)). Thus by the stronger degree bound, \(V\)
is projectively the natural representation of \(B\) up to twist, \(v=w\): we have absolute irreducibility
by the full matrix span, and \(w^2/2-O(w)>xy\).

In this equality case \(B\) matches the ambient root series. In ambient A, the supporting representation
is not self-dual even over the algebraic closure for the finite supporting special groups; in signed
types the form symmetry and dimension determine the series (and type A natural is not self-dual). In
characteristic two a symplectic natural \(B\) cannot preserve the orthogonal quadratic (see also the
field-transfer transitivity test in \Cref{sec:forward-completion}; over its natural finite field an
invariant quadratic, even allowing coefficients
in the closure, would have constant evaluation on nonzero vectors). Such finite-group preservation on
perfect linear lifts would be exact. Now apply \Cref{w:natural-root}: either
\(B=S\) or it comes from a proper root of the ambient Steinberg map, of field height strictly between
zero and the ambient field height \(f\), fixing the supports pointwise projectively. The latter is
impossible by restricted-weight uniqueness: on a supporting group the digit is natural at field exponent
zero (parameter \(p^f\)), not its twist through such a height even up to duality. Frobenius tests on
projective representations here and below agree linearly on universal covers by perfection. This proves
the claim.

\end{proof}

The field-transfer and exchanged-pair overgroup arguments use
\Cref{fa:degree-gap,fa:coefficient-twists,fa:ordered-factor-maximal}.  They complete the remaining
classical exclusions in \Cref{sec:forward-completion}.
